\documentclass[11pt]{article}

% ------------------------------------------------------------
% Packages
% ------------------------------------------------------------
\usepackage[margin=1in]{geometry}
\usepackage{amsmath,amssymb,amsthm}
\usepackage{array}
\usepackage{booktabs}
\usepackage{longtable}
\usepackage{tabularx}
\usepackage{enumitem}
\usepackage{xcolor}
\usepackage{hyperref}
\usepackage{titlesec}
\usepackage{multicol}
\usepackage{fancyhdr}

% ------------------------------------------------------------
% General formatting
% ------------------------------------------------------------
\setlength{\parindent}{0pt}
\setlength{\parskip}{5pt}

\hypersetup{
    colorlinks=true,
    linkcolor=blue,
    urlcolor=blue
}

\pagestyle{fancy}
\fancyhf{}
\lhead{Calculus I for Biology and Social Sciences}
\rhead{Course Structure}
\cfoot{\thepage}

\titleformat{\section}
  {\large\bfseries}
  {\thesection}
  {0.5em}
  {}

\titleformat{\subsection}
  {\normalsize\bfseries}
  {\thesubsection}
  {0.5em}
  {}

% ------------------------------------------------------------
% Custom commands
% ------------------------------------------------------------
\newcommand{\important}{\textbf{Core}}
\newcommand{\optional}{\textit{Optional}}
\newcommand{\bioapp}{\textit{Biological/Application Context}}

% ------------------------------------------------------------
% Document
% ------------------------------------------------------------
\begin{document}

\begin{center}
    {\LARGE\bfseries Calculus I for Biology and Social Sciences}\\[4pt]
    {\large Full Course Structure and Learning Framework}\\[8pt]
    \textit{Based primarily on Claudia Neuhauser and Marcus L. Roper,}\\
    \textit{Calculus for Biology and Medicine, Fourth Edition}\\[8pt]
\end{center}

\hrule
\vspace{0.5cm}

% ============================================================
\section{Course Description}

This course develops the fundamental ideas and techniques of single-variable
calculus with emphasis on applications to biological, health, behavioral,
environmental, and other quantitative sciences.

The course is organized around the central calculus ideas of:

\begin{enumerate}[label=\arabic*.]
    \item functions and mathematical modeling,
    \item limits and continuity,
    \item differentiation and rates of change,
    \item applications of derivatives,
    \item antiderivatives,
    \item definite integrals,
    \item the Fundamental Theorem of Calculus, and
    \item applications of integration.
\end{enumerate}

The course uses the life-sciences orientation of the textbook: mathematical
ideas are introduced through meaningful applications and are subsequently
developed mathematically and applied to models.

The main mathematical progression is

\[
\boxed{
\text{Functions}
\longrightarrow
\text{Limits}
\longrightarrow
\text{Derivatives}
\longrightarrow
\text{Applications}
\longrightarrow
\text{Integrals}
}
\]

The course should emphasize not only symbolic computation but also
interpretation, graphical reasoning, dimensional analysis, model construction,
and communication of quantitative conclusions.

% ============================================================
\section{Textbook Alignment}

The primary textbook is:

\begin{quote}
Claudia Neuhauser and Marcus L. Roper,
\textit{Calculus for Biology and Medicine},
Fourth Edition, Pearson.
\end{quote}

The textbook is explicitly designed around applications in biology and
medicine and emphasizes mathematical modeling, interpretation of mathematical
models, data, and quantitative reasoning.

For a one-semester calculus course with an integration emphasis, the textbook
recommends:

\begin{itemize}
    \item Chapters 3--4,
    \item Sections 5.1--5.6,
    \item Section 5.10,
    \item Sections 6.1--6.3.
\end{itemize}

Chapter 1 is principally prerequisite/precalculus material. Chapter 2 is
useful as an introduction to recurrence models and mathematical modeling,
although it can be compressed or assigned as preparatory work.

Thus the recommended Calculus I pathway is:

\[
\boxed{
\text{Chapter 1 Review}
\rightarrow
\text{Chapter 2 Modeling Bridge}
\rightarrow
\text{Chapter 3}
\rightarrow
\text{Chapter 4}
\rightarrow
\text{Chapter 5}
\rightarrow
\text{Chapter 6}
}
\]

% ============================================================
\section{Overall Course Structure}

\begin{longtable}{p{0.8cm}p{3.0cm}p{4.0cm}p{5.2cm}}
\toprule
\textbf{Unit} & \textbf{Textbook} & \textbf{Major Theme}
& \textbf{Central Concepts} \\
\midrule
1 & Ch. 1 & Prerequisite mathematics
& Algebra, functions, trigonometry, exponentials, logarithms, graphing \\

2 & Ch. 2 & Mathematical modeling bridge
& Discrete growth, recurrence equations, sequences, limits, model assumptions \\

3 & Ch. 3 & Limits and continuity
& Limits, continuity, limits at infinity, trigonometric limits, IVT \\

4 & Ch. 4 & Differentiation
& Derivative, rates of change, differentiation rules, chain rule, implicit differentiation \\

5 & Ch. 5 & Applications of differentiation
& Extrema, MVT, monotonicity, concavity, optimization, graphing, antiderivatives \\

6 & Ch. 6 & Integration
& Definite integral, Riemann sums, FTC, cumulative change, average value \\

\bottomrule
\end{longtable}

% ============================================================
\section{Prerequisite Knowledge}

Students entering Calculus I should be able to work comfortably with the
following topics.

\subsection{Algebra}

\begin{itemize}
    \item Real numbers and intervals.
    \item Algebraic manipulation.
    \item Factoring.
    \item Solving linear and quadratic equations.
    \item Solving inequalities.
    \item Fractions and rational expressions.
    \item Exponents and radicals.
    \item Scientific notation.
    \item Coordinate geometry.
    \item Equations of lines.
    \item Basic polynomial manipulation.
\end{itemize}

\subsection{Functions}

Students should understand:

\begin{itemize}
    \item function notation,
    \item domain and range,
    \item composition,
    \item inverse functions,
    \item polynomial functions,
    \item rational functions,
    \item power functions,
    \item exponential functions,
    \item logarithmic functions, and
    \item basic transformations of graphs.
\end{itemize}

\subsection{Trigonometry}

Students should know:

\begin{itemize}
    \item radians and degrees,
    \item basic trigonometric functions,
    \item unit-circle values,
    \item basic trigonometric identities,
    \item graphs of sine, cosine, and tangent,
    \item inverse trigonometric functions.
\end{itemize}

\subsection{Graphing and Data Interpretation}

Students should be able to:

\begin{itemize}
    \item interpret a graph,
    \item identify increasing and decreasing behavior,
    \item estimate slopes,
    \item interpret intercepts,
    \item recognize transformations,
    \item read logarithmic and semilogarithmic graphs,
    \item translate a verbal description into a qualitative graph.
\end{itemize}

% ============================================================
\section{Course-Level Learning Objectives}

By the end of Calculus I, students should be able to:

\begin{enumerate}[label=\textbf{LO\arabic*.}]
    \item Analyze and manipulate elementary functions used in quantitative
          science.
          
    \item Interpret graphs and mathematical relationships in biological and
          social-science contexts.
          
    \item Construct simple mathematical models from verbal descriptions.
    
    \item Explain the intuitive meaning of a limit as a quantity approaches
          a specified value.
          
    \item Evaluate limits algebraically, graphically, and numerically.
    
    \item Determine whether a function is continuous at a point or on an
          interval.
          
    \item Explain the derivative as an instantaneous rate of change and as
          the slope of a tangent line.
          
    \item Compute derivatives using the fundamental differentiation rules.
    
    \item Apply the product, quotient, and chain rules.
    
    \item Differentiate implicit relationships and solve related-rate
          problems.
          
    \item Interpret higher derivatives in terms of changes in rates.
    
    \item Use derivatives to identify extrema, monotonicity, concavity, and
          inflection points.
          
    \item Apply the Mean Value Theorem and Extreme Value Theorem conceptually
          and computationally.
          
    \item Formulate and solve optimization problems arising from scientific
          models.
          
    \item Analyze and sketch functions using derivative information.
    
    \item Explain the meaning of an antiderivative.
    
    \item Interpret a definite integral as accumulated change and as area
          under a curve.
          
    \item Approximate definite integrals using Riemann sums.
    
    \item Apply the Fundamental Theorem of Calculus.
    
    \item Calculate cumulative change and average values using definite
          integrals.
          
    \item Translate between graphical, numerical, symbolic, and verbal
          representations of calculus concepts.
          
    \item Select an appropriate calculus method for a scientific problem and
          justify the choice.
          
    \item Interpret the mathematical result in the original scientific
          context, including units and limitations of the model.
\end{enumerate}

% ============================================================
\section{Module 1: Precalculus and Mathematical Foundations}

\subsection{Textbook Coverage}

\textbf{Chapter 1: Preview and Review}

\begin{itemize}
    \item 1.1 Precalculus Skills Diagnostic Test
    \item 1.2 Preliminaries
    \item 1.3 Elementary Functions
    \item 1.4 Graphing
\end{itemize}

\subsection{Major Concepts}

\begin{itemize}
    \item Real numbers.
    \item Lines in the plane.
    \item Equations of circles.
    \item Trigonometry.
    \item Exponential and logarithmic functions.
    \item Complex numbers and quadratic equations.
    \item Function notation.
    \item Polynomial and rational functions.
    \item Power functions.
    \item Inverse functions.
    \item Function transformations.
    \item Logarithmic and semilogarithmic scales.
    \item Graphical representation of scientific data.
\end{itemize}

\subsection{Learning Objectives}

Students should be able to:

\begin{enumerate}
    \item manipulate algebraic expressions accurately;
    \item solve equations and inequalities;
    \item evaluate and compose functions;
    \item determine domains and ranges;
    \item work with exponential and logarithmic functions;
    \item use basic trigonometric identities;
    \item transform graphs;
    \item interpret logarithmic scales;
    \item convert scientific descriptions into graphs.
\end{enumerate}

\subsection{Prerequisites}

This module is itself the prerequisite foundation for the calculus course.

\subsection{Potentially Difficult Topics}

\begin{itemize}
    \item Function composition and inverse functions.
    \item Exponential versus logarithmic representations.
    \item Logarithmic graph transformations.
    \item Trigonometric identities.
    \item Translating verbal descriptions into graphs.
\end{itemize}

\subsection{Recommended Instructional Role}

Most of Chapter 1 should be treated as prerequisite review rather than as a
full calculus module. Students can complete a diagnostic assessment and
receive targeted remediation.

% ============================================================
\section{Module 2: Discrete Models and Mathematical Modeling}

\subsection{Textbook Coverage}

\textbf{Chapter 2: Discrete-Time Models, Sequences, and Difference Equations}

\begin{itemize}
    \item 2.1 Exponential Growth and Decay
    \item 2.2 Sequences
    \item 2.3 Modeling with Recurrence Equations
\end{itemize}

\subsection{Major Concepts}

\begin{itemize}
    \item Discrete-time population growth.
    \item Exponential growth and decay.
    \item Recurrence equations.
    \item Sequences.
    \item Limits of sequences.
    \item Sigma notation.
    \item Density-dependent population growth.
    \item Beverton--Holt model.
    \item Discrete logistic equation.
    \item Model assumptions.
    \item Comparing mathematical models with data.
\end{itemize}

\subsection{Learning Objectives}

Students should be able to:

\begin{enumerate}
    \item distinguish discrete-time from continuous-time models;
    \item formulate simple recurrence equations;
    \item calculate terms of a sequence;
    \item interpret long-term behavior of sequences;
    \item use sigma notation;
    \item identify assumptions in a mathematical model;
    \item translate a verbal scientific description into a recurrence model;
    \item interpret model parameters;
    \item compare a model with observed data.
\end{enumerate}

\subsection{Prerequisites}

\begin{itemize}
    \item Algebra.
    \item Functions.
    \item Exponential functions.
    \item Graph interpretation.
\end{itemize}

\subsection{Potentially Difficult Topics}

\begin{itemize}
    \item Translating verbal assumptions into equations.
    \item Understanding recurrence equations.
    \item Distinguishing model behavior from actual biological behavior.
    \item Interpreting parameters rather than simply calculating them.
\end{itemize}

\subsection{Course Role}

This module serves as a bridge between prerequisite mathematics and
continuous-time calculus. It establishes the modeling perspective that should
continue throughout the course.

% ============================================================
\section{Module 3: Limits and Continuity}

\subsection{Textbook Coverage}

\textbf{Chapter 3: Limits and Continuity}

\begin{itemize}
    \item 3.1 Limits
    \item 3.2 Continuity
    \item 3.3 Limits at Infinity
    \item 3.4 Trigonometric Limits and the Sandwich Theorem
    \item 3.5 Properties of Continuous Functions
    \item 3.6 Formal Definition of Limits \optional
\end{itemize}

\subsection{Major Concepts}

\begin{itemize}
    \item Informal definition of a limit.
    \item One-sided behavior.
    \item Limit laws.
    \item Algebraic techniques for evaluating limits.
    \item Indeterminate forms.
    \item Continuity.
    \item Limits at infinity.
    \item Infinite limits and asymptotic behavior.
    \item Trigonometric limits.
    \item Sandwich Theorem.
    \item Intermediate Value Theorem.
    \item Bisection method.
    \item Formal $\epsilon$--$\delta$ definition of a limit.
\end{itemize}

\subsection{Learning Objectives}

Students should be able to:

\begin{enumerate}
    \item explain the meaning of
          \[
          \lim_{x\to a}f(x);
          \]
    \item estimate limits graphically and numerically;
    \item evaluate limits using algebraic limit laws;
    \item recognize when direct substitution is insufficient;
    \item evaluate limits involving trigonometric functions;
    \item analyze end behavior;
    \item determine continuity;
    \item apply the Intermediate Value Theorem;
    \item explain how the bisection method locates roots;
    \item connect limits to the later concept of derivatives.
\end{enumerate}

\subsection{Prerequisites}

\begin{itemize}
    \item Functions and function notation.
    \item Algebraic manipulation.
    \item Factoring.
    \item Rational expressions.
    \item Graph interpretation.
    \item Trigonometry.
\end{itemize}

\subsection{Difficult Topics}

\begin{itemize}
    \item Conceptual meaning of a limit.
    \item Difference between $f(a)$ and $\lim_{x\to a}f(x)$.
    \item Algebraic manipulation of indeterminate forms.
    \item Limits at infinity.
    \item Trigonometric limits.
    \item Formal definition of a limit.
\end{itemize}

\subsection{Essential Conceptual Transition}

The key conceptual transition is:

\[
\text{average rate of change}
\quad\longrightarrow\quad
\text{instantaneous rate of change}.
\]

This prepares students for the derivative.

% ============================================================
\section{Module 4: Differentiation}

\subsection{Textbook Coverage}

\textbf{Chapter 4: Differentiation}

\begin{itemize}
    \item 4.1 Formal Definition of the Derivative
    \item 4.2 Properties of the Derivative
    \item 4.3 Power Rule and Basic Rules
    \item 4.4 Product and Quotient Rules
    \item 4.5 Chain Rule
    \item 4.6 Implicit Functions and Implicit Differentiation
    \item 4.7 Higher Derivatives
    \item 4.8 Derivatives of Trigonometric Functions
    \item 4.9 Derivatives of Exponential Functions
    \item 4.10 Derivatives of Inverse and Logarithmic Functions
    \item 4.11 Linear Approximation and Error Propagation
\end{itemize}

\subsection{Major Concepts}

\begin{itemize}
    \item Difference quotient.
    \item Derivative as a limit.
    \item Tangent-line slope.
    \item Instantaneous rate of change.
    \item Differentiability.
    \item Relationship between continuity and differentiability.
    \item Power rule.
    \item Constant and sum rules.
    \item Product rule.
    \item Quotient rule.
    \item Chain rule.
    \item Implicit differentiation.
    \item Related rates.
    \item Higher derivatives.
    \item Trigonometric derivatives.
    \item Exponential and logarithmic derivatives.
    \item Linear approximation.
    \item Error propagation.
\end{itemize}

\subsection{Learning Objectives}

Students should be able to:

\begin{enumerate}
    \item define the derivative using a limit;
    \item interpret derivatives geometrically;
    \item interpret derivatives as rates of change;
    \item determine whether a function is differentiable;
    \item compute derivatives using standard rules;
    \item combine multiple differentiation rules;
    \item apply the chain rule to composite functions;
    \item differentiate implicitly defined functions;
    \item solve related-rate problems;
    \item calculate and interpret higher derivatives;
    \item differentiate trigonometric functions;
    \item differentiate exponential and logarithmic functions;
    \item construct linear approximations;
    \item estimate propagation of measurement errors.
\end{enumerate}

\subsection{Prerequisites}

\begin{itemize}
    \item Chapter 3 limits.
    \item Algebra.
    \item Function composition.
    \item Trigonometry.
    \item Exponential and logarithmic functions.
\end{itemize}

\subsection{Difficult Topics}

\begin{enumerate}
    \item Understanding the derivative conceptually rather than as a
          differentiation algorithm.
    \item Connecting the derivative to a physical or biological rate.
    \item The chain rule.
    \item Nested chain-rule problems.
    \item Implicit differentiation.
    \item Related rates.
    \item Selecting the correct combination of derivative rules.
\end{enumerate}

\subsection{Central Mathematical Relationship}

The derivative is defined by

\[
f'(x)
=
\lim_{h\to 0}
\frac{f(x+h)-f(x)}{h}.
\]

It represents both

\[
\boxed{\text{slope of tangent line}}
\]

and

\[
\boxed{\text{instantaneous rate of change}}.
\]

For a biological quantity $N(t)$,

\[
N'(t)=\frac{dN}{dt}
\]

represents the instantaneous rate at which the population changes.

% ============================================================
\section{Module 5: Applications of Differentiation}

\subsection{Textbook Coverage}

\textbf{Chapter 5: Applications of Differentiation}

\begin{itemize}
    \item 5.1 Extrema and the Mean Value Theorem
    \item 5.2 Monotonicity and Concavity
    \item 5.3 Extrema and Inflection Points
    \item 5.4 Optimization
    \item 5.5 L'H\^opital's Rule
    \item 5.6 Graphing and Asymptotes
    \item 5.10 Antiderivatives
\end{itemize}

The following sections can be treated as extensions:

\begin{itemize}
    \item 5.7 Recurrence Equations: Stability
    \item 5.8 Newton--Raphson Method
    \item 5.9 Modeling Biological Systems Using Differential Equations
\end{itemize}

\subsection{Major Concepts}

\begin{itemize}
    \item Absolute and local extrema.
    \item Extreme Value Theorem.
    \item Mean Value Theorem.
    \item Critical points.
    \item First derivative and monotonicity.
    \item Second derivative and concavity.
    \item Inflection points.
    \item Optimization.
    \item Indeterminate forms.
    \item L'H\^opital's Rule.
    \item Asymptotes.
    \item Curve sketching.
    \item Antiderivatives.
\end{itemize}

\subsection{Learning Objectives}

Students should be able to:

\begin{enumerate}
    \item identify critical points;
    \item distinguish local and global extrema;
    \item apply the Extreme Value Theorem;
    \item apply the Mean Value Theorem;
    \item determine intervals of increase and decrease;
    \item determine concavity;
    \item identify inflection points;
    \item use first- and second-derivative information to classify extrema;
    \item solve scientific optimization problems;
    \item interpret an optimal solution in context;
    \item analyze asymptotic behavior;
    \item construct a qualitative graph using calculus;
    \item use L'H\^opital's Rule for appropriate indeterminate limits;
    \item recognize antiderivatives as inverse operations to differentiation.
\end{enumerate}

\subsection{Prerequisites}

\begin{itemize}
    \item Derivatives.
    \item Chain rule.
    \item Algebra.
    \item Graph interpretation.
    \item Limits.
\end{itemize}

\subsection{Difficult Topics}

\begin{itemize}
    \item Distinguishing local from absolute extrema.
    \item Selecting critical points correctly.
    \item Applying the Mean Value Theorem.
    \item Relating the first derivative to function behavior.
    \item Relating the second derivative to concavity.
    \item Translating word problems into optimization equations.
    \item Interpreting an optimization result scientifically.
    \item Combining multiple derivative tests in curve sketching.
\end{itemize}

\subsection{Optimization Framework}

A general optimization problem should follow:

\[
\text{Scientific situation}
\rightarrow
\text{Variables}
\rightarrow
\text{Objective function}
\rightarrow
\text{Constraints}
\rightarrow
\text{One-variable function}
\rightarrow
\text{Derivative}
\rightarrow
\text{Critical points}
\rightarrow
\text{Optimal solution}
\rightarrow
\text{Interpretation}.
\]

% ============================================================
\section{Module 6: Antiderivatives and Integration}

\subsection{Textbook Coverage}

\textbf{Chapter 5.10 and Chapter 6}

\begin{itemize}
    \item 5.10 Antiderivatives
    \item 6.1 The Definite Integral
    \item 6.2 The Fundamental Theorem of Calculus
    \item 6.3 Applications of Integration
\end{itemize}

\subsection{Major Concepts}

\begin{itemize}
    \item Antiderivatives.
    \item Indefinite integrals.
    \item Accumulation.
    \item Area under a curve.
    \item Riemann sums.
    \item Definite integrals.
    \item Properties of the integral.
    \item Fundamental Theorem of Calculus, Part I.
    \item Fundamental Theorem of Calculus, Part II.
    \item Cumulative change.
    \item Average value.
    \item Area between curves \optional.
    \item Volume \optional.
    \item Arc length \optional.
\end{itemize}

\subsection{Learning Objectives}

Students should be able to:

\begin{enumerate}
    \item explain the relationship between derivatives and antiderivatives;
    \item find basic antiderivatives;
    \item interpret a definite integral as accumulated change;
    \item approximate an integral using Riemann sums;
    \item distinguish a definite integral from an indefinite integral;
    \item evaluate definite integrals using antiderivatives;
    \item apply both parts of the Fundamental Theorem of Calculus;
    \item calculate cumulative change from a rate;
    \item calculate average values of functions;
    \item interpret an integral in a biological or scientific context;
    \item determine appropriate units for an integral.
\end{enumerate}

\subsection{Prerequisites}

\begin{itemize}
    \item Limits.
    \item Derivatives.
    \item Antiderivatives.
    \item Algebra.
    \item Function interpretation.
\end{itemize}

\subsection{Difficult Topics}

The most important conceptual difficulties are:

\begin{enumerate}
    \item Understanding the definite integral as accumulation rather than
          merely ``area.''
    \item Understanding Riemann sums.
    \item Connecting a rate function to accumulated quantity.
    \item Distinguishing $\int f(x)\,dx$ from
          $\int_a^b f(x)\,dx$.
    \item Understanding both directions of the Fundamental Theorem of
          Calculus.
    \item Setting up integrals from scientific descriptions.
\end{enumerate}

\subsection{Central Relationship}

The Fundamental Theorem of Calculus connects differentiation and integration:

\[
\boxed{
\frac{d}{dx}
\left(
\int_a^x f(t)\,dt
\right)
=
f(x)
}
\]

and, when $F'(x)=f(x)$,

\[
\boxed{
\int_a^b f(x)\,dx
=
F(b)-F(a).
}
\]

Thus,

\[
\boxed{
\text{differentiation}
\longleftrightarrow
\text{integration}
}
\]

is the central structural idea of the second half of Calculus I.

% ============================================================
\section{Module 7: Integration Applications for Life and Social Sciences}

\subsection{Core Applications}

The course should prioritize applications that are directly meaningful for
scientific students.

\subsubsection{Cumulative Change}

If $R(t)$ is a rate of change, then

\[
\Delta Q
=
\int_a^b R(t)\,dt.
\]

For a population,

\[
N(b)-N(a)
=
\int_a^b N'(t)\,dt.
\]

For drug concentration or flow-rate models, the same principle can be used
to recover cumulative quantities.

\subsubsection{Average Value}

For a function $f$ on $[a,b]$,

\[
f_{\mathrm{avg}}
=
\frac{1}{b-a}
\int_a^b f(x)\,dx.
\]

Students should be able to distinguish an instantaneous value $f(x)$ from
the average value over an interval.

\subsubsection{Area}

The definite integral can represent accumulated area:

\[
A
=
\int_a^b f(x)\,dx
\]

when $f(x)\geq 0$.

More generally, the signed area between two functions is

\[
A
=
\int_a^b
\left[f(x)-g(x)\right]dx.
\]

\subsection{Learning Objectives}

Students should be able to:

\begin{enumerate}
    \item identify a rate from a scientific description;
    \item construct an integral representing accumulated change;
    \item evaluate the integral;
    \item determine the correct units;
    \item interpret the result in context;
    \item calculate and interpret average values;
    \item distinguish net accumulation from geometric area.
\end{enumerate}

% ============================================================
\section{Prerequisite Dependency Structure}

The conceptual dependency structure of the course is:

\[
\begin{array}{c}
\text{Algebra, Functions, Trigonometry} \\[3pt]
\downarrow \\[3pt]
\text{Graphing and Modeling} \\[3pt]
\downarrow \\[3pt]
\text{Limits and Continuity} \\[3pt]
\downarrow \\[3pt]
\text{Derivative} \\[3pt]
\downarrow \\[3pt]
\text{Differentiation Rules} \\[3pt]
\downarrow \\[3pt]
\text{Applications of Derivatives} \\[3pt]
\downarrow \\[3pt]
\text{Antiderivatives} \\[3pt]
\downarrow \\[3pt]
\text{Definite Integral} \\[3pt]
\downarrow \\[3pt]
\text{Fundamental Theorem of Calculus} \\[3pt]
\downarrow \\[3pt]
\text{Applications of Integration}
\end{array}
\]

A more detailed dependency graph is:

\[
\boxed{
\begin{array}{ccccc}
\text{Algebra} & \rightarrow & \text{Functions} &
\rightarrow & \text{Limits} \\[3pt]
&&&& \downarrow \\[3pt]
\text{Trigonometry} & \rightarrow & \text{Exponential/Logarithmic Functions}
&& \text{Derivative} \\[3pt]
&&&& \downarrow \\[3pt]
&&&& \text{Differentiation Rules} \\[3pt]
&&&& \downarrow \\[3pt]
&&&& \text{Optimization/Graphing} \\[3pt]
&&&& \downarrow \\[3pt]
&&&& \text{Antiderivatives} \\[3pt]
&&&& \downarrow \\[3pt]
&&&& \text{Definite Integrals} \\[3pt]
&&&& \downarrow \\[3pt]
&&&& \text{FTC and Applications}
\end{array}
}
\]

% ============================================================
\section{Major Difficult Topics}

The following topics should receive additional instructional time and
practice.

\begin{longtable}{p{0.7cm}p{4.0cm}p{6.5cm}p{3.0cm}}
\toprule
\textbf{\#} & \textbf{Topic} & \textbf{Why It Is Difficult}
& \textbf{Recommended Support} \\
\midrule

1 & Limits
& Students must transition from numerical substitution to a conceptual
notion of approaching a value.
& Graphical and numerical examples \\

2 & Derivative as a limit
& Requires students to connect average and instantaneous rates.
& Secant-to-tangent visualizations \\

3 & Chain Rule
& Requires identification of nested functions and multiple differentiation
steps.
& Layered function notation \\

4 & Implicit Differentiation
& The dependent variable is not isolated.
& Step-by-step symbolic practice \\

5 & Related Rates
& Students must construct the mathematical relationship before differentiating.
& Modeling diagrams and units \\

6 & Optimization
& Requires translation from words to mathematical structure.
& Explicit problem-solving framework \\

7 & Concavity and Inflection
& Students must interpret the second derivative conceptually.
& Graph/derivative correspondence \\

8 & Riemann Sums
& Requires a transition from finite sums to limiting processes.
& Numerical approximations and graphical areas \\

9 & Fundamental Theorem of Calculus
& Students must connect two apparently different operations:
differentiation and integration.
& Multiple representations \\

10 & Integral Setup
& The challenge is often modeling rather than integration itself.
& Rate-to-accumulation problems \\

\bottomrule
\end{longtable}

% ============================================================
\section{Suggested 15-Week Course Schedule}

\begin{longtable}{p{0.7cm}p{3.2cm}p{5.5cm}p{4.2cm}}
\toprule
\textbf{Week} & \textbf{Module} & \textbf{Topics}
& \textbf{Major Learning Goal} \\
\midrule

1
& Foundations
& Chapter 1 diagnostic; algebra; functions; graphing
& Establish prerequisite fluency \\

2
& Modeling Bridge
& Exponential growth; recurrence equations; sequences; modeling
& Translate scientific descriptions into mathematical models \\

3
& Limits
& Limit concept; numerical and graphical limits; limit laws
& Understand and calculate limits \\

4
& Limits and Continuity
& Limits at infinity; trigonometric limits; continuity; IVT
& Connect limits to continuity \\

5
& Derivatives
& Derivative definition; tangent lines; rates of change
& Understand derivative conceptually \\

6
& Differentiation
& Power rule; product rule; quotient rule
& Compute derivatives efficiently \\

7
& Differentiation
& Chain rule; trigonometric; exponential and logarithmic derivatives
& Differentiate composite functions \\

8
& Differentiation Applications
& Implicit differentiation; related rates; higher derivatives
& Model changing quantities \\

9
& Extrema
& Critical points; EVT; MVT; monotonicity
& Analyze function behavior \\

10
& Concavity and Optimization
& Concavity; inflection points; optimization
& Solve scientific optimization problems \\

11
& Graphing and Antiderivatives
& Curve sketching; asymptotes; antiderivatives
& Connect derivatives and antiderivatives \\

12
& Integration
& Area problem; Riemann sums; definite integral
& Understand integration conceptually \\

13
& Fundamental Theorem
& FTC Part I and Part II; indefinite integrals
& Connect differentiation and integration \\

14
& Integration Applications
& Cumulative change; average value; area
& Interpret integrals scientifically \\

15
& Synthesis
& Integrated modeling problems; review; applications
& Select and justify appropriate calculus methods \\

\bottomrule
\end{longtable}

% ============================================================
\section{Detailed Weekly Learning Sequence}

\subsection*{Week 1: Prerequisite Readiness}

\textbf{Concepts:}
\begin{itemize}
    \item Functions
    \item Algebraic manipulation
    \item Graphs
    \item Exponents and logarithms
    \item Trigonometry
\end{itemize}

\textbf{Assessment:}
Precalculus diagnostic and prerequisite skills quiz.

\subsection*{Week 2: Mathematical Modeling}

\textbf{Concepts:}
\begin{itemize}
    \item Exponential growth and decay
    \item Recurrence equations
    \item Sequences
    \item Model assumptions
\end{itemize}

\textbf{Application:}
Population growth and physiological/drug models.

\subsection*{Weeks 3--4: Limits and Continuity}

\textbf{Core question:}

\begin{quote}
What does it mean for a quantity to approach a value?
\end{quote}

Students move from intuitive graphical ideas to algebraic limit laws and
continuity.

\subsection*{Weeks 5--7: Differentiation}

\textbf{Core question:}

\begin{quote}
How quickly is a quantity changing at an instant?
\end{quote}

The progression should be

\[
\text{average rate}
\rightarrow
\text{limit}
\rightarrow
\text{instantaneous rate}
\rightarrow
\text{derivative}
\rightarrow
\text{differentiation rules}.
\]

\subsection*{Week 8: Implicit Differentiation and Related Rates}

Students apply derivatives to relationships in which the dependent variable
is not explicitly isolated and to systems whose quantities change
simultaneously.

\subsection*{Weeks 9--10: Applications of Differentiation}

Students use derivatives to answer questions such as:

\begin{itemize}
    \item When is a population increasing most rapidly?
    \item When does a quantity reach a maximum?
    \item When is a function increasing or decreasing?
    \item When does growth accelerate or decelerate?
    \item What parameter value optimizes a biological process?
\end{itemize}

\subsection*{Week 11: Antiderivatives}

Students reverse differentiation:

\[
f'(x)=g(x)
\qquad\Longrightarrow\qquad
f(x)=\int g(x)\,dx.
\]

\subsection*{Weeks 12--13: Definite Integrals and FTC}

Students progress from

\[
\text{rectangles}
\rightarrow
\text{Riemann sums}
\rightarrow
\text{definite integral}
\rightarrow
\text{Fundamental Theorem}.
\]

\subsection*{Weeks 14--15: Applications and Synthesis}

The final part of the course should emphasize interpretation and modeling,
not merely symbolic integration.

% ============================================================
\section{Core Concept Inventory}

\begin{longtable}{p{4.0cm}p{8.8cm}}
\toprule
\textbf{Concept Category} & \textbf{Concepts Students Should Master} \\
\midrule

Functions
& Domain, range, composition, inverse functions, transformations,
polynomial, rational, power, exponential, logarithmic, trigonometric functions \\

Modeling
& Variables, parameters, assumptions, recurrence equations, model
interpretation, comparison with data \\

Limits
& Limit notation, one-sided limits, limit laws, infinite limits, limits at
infinity, trigonometric limits \\

Continuity
& Continuity at a point, continuity on an interval, Intermediate Value
Theorem \\

Derivatives
& Difference quotient, derivative, tangent line, instantaneous rate,
differentiability \\

Differentiation
& Power, product, quotient, chain, implicit differentiation, related rates,
higher derivatives \\

Derivative Applications
& Critical points, extrema, MVT, monotonicity, concavity, inflection points,
optimization, asymptotes, graphing \\

Antiderivatives
& Antiderivatives, indefinite integrals, constant of integration \\

Integration
& Area, Riemann sums, definite integrals, accumulation \\

Fundamental Theorem
& FTC Part I, FTC Part II, derivative-integral relationship \\

Integration Applications
& Cumulative change, average value, area, scientific interpretation \\

\bottomrule
\end{longtable}

% ============================================================
\section{Mathematical Modeling Competencies}

Because the textbook places unusually strong emphasis on modeling, modeling
should be treated as a course-wide competency rather than as a single
chapter.

Students should learn the following modeling cycle:

\[
\boxed{
\begin{array}{c}
\text{Real-world phenomenon}\\
\downarrow\\
\text{Identify quantities and variables}\\
\downarrow\\
\text{State assumptions}\\
\downarrow\\
\text{Construct mathematical model}\\
\downarrow\\
\text{Analyze mathematically}\\
\downarrow\\
\text{Compare with data or expected behavior}\\
\downarrow\\
\text{Interpret results}\\
\downarrow\\
\text{Evaluate limitations}
\end{array}
}
\]

For every major calculus topic, at least one problem should require students
to complete this cycle.

% ============================================================
\section{Application Domains}

The textbook provides extensive life-science applications. For a broader
Biology and Social Sciences course, the same calculus structures can be
organized around the following application domains.

\subsection{Biology}

\begin{itemize}
    \item Population growth.
    \item Cell growth.
    \item Population dynamics.
    \item Species-area relationships.
    \item Physiological rates.
    \item Drug concentration and metabolism.
    \item Chemical reaction rates.
    \item Ecological processes.
    \item Epidemiological processes.
\end{itemize}

\subsection{Health and Medicine}

\begin{itemize}
    \item Drug absorption and elimination.
    \item Concentration and dosage models.
    \item Physiological rates.
    \item Cumulative exposure.
    \item Medical optimization problems.
\end{itemize}

\subsection{Behavioral and Social Sciences}

For a Biology and Social Sciences version of the course, analogous calculus
problems can be constructed around:

\begin{itemize}
    \item population and demographic growth;
    \item migration;
    \item adoption of information;
    \item social contagion;
    \item behavioral change;
    \item resource allocation;
    \item marginal effects;
    \item cost-benefit optimization;
    \item accumulation of social or economic quantities.
\end{itemize}

These examples should preserve the mathematical structure of the textbook
while broadening the application context.

% ============================================================
\section{Assessment Structure}

A recommended assessment system is:

\begin{center}
\begin{tabular}{lc}
\toprule
\textbf{Assessment Component} & \textbf{Weight} \\
\midrule
Prerequisite/Diagnostic Assessment & 5\% \\
Weekly Homework and Practice & 15\% \\
Conceptual Quizzes & 10\% \\
Application/Modeling Assignments & 15\% \\
Midterm Examination I & 15\% \\
Midterm Examination II & 15\% \\
Final Examination & 25\% \\
\midrule
\textbf{Total} & \textbf{100\%} \\
\bottomrule
\end{tabular}
\end{center}

The exact percentages can be modified according to institutional policy.

% ============================================================
\section{Assessment by Learning Domain}

\begin{longtable}{p{4cm}p{7cm}p{3cm}}
\toprule
\textbf{Domain} & \textbf{Student Performance}
& \textbf{Assessment} \\
\midrule

Procedural Fluency
& Compute limits, derivatives, antiderivatives, and integrals
& Quizzes/Homework \\

Conceptual Understanding
& Explain limits, derivatives, integrals, and FTC
& Conceptual Questions \\

Graphical Reasoning
& Interpret functions and derivative/integral graphs
& Quizzes/Exams \\

Modeling
& Translate scientific situations into equations
& Modeling Assignments \\

Interpretation
& Explain results in scientific language
& Application Problems \\

Problem Solving
& Select and justify an appropriate calculus method
& Exams/Projects \\

Communication
& Clearly present mathematical reasoning and units
& Written Assignments \\

\bottomrule
\end{longtable}

% ============================================================
\section{Recommended Teaching Strategy}

The course should use four representations repeatedly:

\[
\boxed{
\text{Symbolic}
\quad
\text{Graphical}
\quad
\text{Numerical}
\quad
\text{Verbal}
}
\]

For example, when introducing a derivative:

\begin{itemize}
    \item \textbf{Verbal:} instantaneous rate of change;
    \item \textbf{Graphical:} slope of a tangent line;
    \item \textbf{Numerical:} limit of secant slopes;
    \item \textbf{Symbolic:}
          \[
          f'(x)=\lim_{h\to0}
          \frac{f(x+h)-f(x)}{h}.
          \]
\end{itemize}

The same four-representation strategy should be used for limits,
differentiation, optimization, and integration.

% ============================================================
\section{Technology Integration}

Technology should support conceptual understanding rather than replace
mathematical reasoning.

Recommended uses include:

\begin{itemize}
    \item graphing calculators;
    \item spreadsheet calculations;
    \item numerical approximation;
    \item graphical exploration;
    \item parameter exploration;
    \item model comparison;
    \item numerical root-finding;
    \item numerical integration.
\end{itemize}

Students should still be required to explain:

\begin{enumerate}
    \item what calculation they performed;
    \item why the method is appropriate;
    \item what the result means;
    \item what the units are;
    \item what assumptions underlie the result.
\end{enumerate}

% ============================================================
\section{Recommended Core versus Optional Material}

\subsection{Core Calculus I}

\begin{itemize}
    \item Chapter 1 prerequisite review.
    \item Chapter 2 modeling bridge, especially Section 2.3.
    \item Chapter 3.
    \item Chapter 4.
    \item Sections 5.1--5.6.
    \item Section 5.10.
    \item Sections 6.1--6.3.
\end{itemize}

\subsection{Optional Enrichment}

Depending on available time:

\begin{itemize}
    \item Section 3.6: Formal definition of limits.
    \item Section 5.7: Stability of recurrence equations.
    \item Section 5.8: Newton--Raphson method.
    \item Section 5.9: Biological differential-equation modeling.
    \item Sections 6.3.4--6.3.6: additional geometric applications.
    \item Section 7.1: substitution rule.
    \item Section 7.5: numerical integration.
\end{itemize}

The textbook itself identifies a number of sections with an asterisk as
optional, including more formal treatments and additional modeling
applications. These can be selected according to course time and student
needs.

% ============================================================
\section{Minimum Competency Standards}

A student successfully completing Calculus I should be able to solve the
following categories of problems without relying on memorized procedures
alone.

\subsection*{Competency 1: Function Analysis}

Given

\[
f(x)=\frac{x^2+1}{x-2},
\]

the student can determine the domain, identify relevant graph behavior, and
interpret the function.

\subsection*{Competency 2: Limit}

Given

\[
\lim_{x\to a}f(x),
\]

the student can estimate and calculate the limit and explain its meaning.

\subsection*{Competency 3: Derivative}

Given a model $N(t)$, the student can calculate

\[
N'(t)
\]

and explain what the derivative means in terms of the modeled quantity.

\subsection*{Competency 4: Optimization}

Given a scientific objective and constraints, the student can construct an
objective function, differentiate it, identify candidate extrema, and
interpret the optimal value.

\subsection*{Competency 5: Accumulation}

Given a rate $R(t)$, the student can construct

\[
\int_a^b R(t)\,dt
\]

and explain the accumulated quantity represented by the integral.

\subsection*{Competency 6: Fundamental Theorem}

The student can move between

\[
f(t)
\quad\leftrightarrow\quad
\int_a^t f(s)\,ds
\quad\leftrightarrow\quad
\frac{d}{dt}
\left[
\int_a^t f(s)\,ds
\right].
\]

% ============================================================
\section{Course Capstone}

The final course assessment should include a small mathematical-modeling
problem in which students must use several calculus ideas.

A recommended structure is:

\begin{enumerate}
    \item Define a real-world quantity $Q(t)$.
    \item Identify the relevant variables and parameters.
    \item State assumptions.
    \item Construct a mathematical model.
    \item Analyze the model using derivatives.
    \item Determine an important maximum, minimum, or rate.
    \item Use integration to calculate cumulative change where appropriate.
    \item Interpret the result.
    \item Discuss model limitations.
\end{enumerate}

The final product should demonstrate that students understand calculus as a
tool for analyzing changing quantities rather than as a collection of
isolated computational techniques.

% ============================================================
\section{Summary Course Map}

\begin{center}
\renewcommand{\arraystretch}{1.5}
\begin{tabular}{c|c|c}
\toprule
\textbf{Stage} & \textbf{Mathematical Question} & \textbf{Core Tool} \\
\midrule
1 & How are quantities related? & Functions \\
2 & What happens as a quantity approaches a value? & Limits \\
3 & Is the function behaving continuously? & Continuity \\
4 & How quickly is a quantity changing? & Derivative \\
5 & When does the quantity increase/decrease fastest? & Derivative Analysis \\
6 & What value optimizes the system? & Optimization \\
7 & What accumulated quantity results from a rate? & Integral \\
8 & How are rates and accumulated quantities related? & FTC \\
9 & How can calculus explain a scientific phenomenon? & Mathematical Modeling \\
\bottomrule
\end{tabular}
\end{center}

\vspace{0.5cm}

The overall intellectual progression of the course is therefore:

\[
\boxed{
\begin{aligned}
\text{Functions}
&\rightarrow
\text{Modeling}
\rightarrow
\text{Limits}
\rightarrow
\text{Continuity}
\\
&\rightarrow
\text{Derivatives}
\rightarrow
\text{Applications}
\rightarrow
\text{Antiderivatives}
\\
&\rightarrow
\text{Definite Integrals}
\rightarrow
\text{Fundamental Theorem}
\rightarrow
\text{Accumulation and Modeling}.
\end{aligned}
}
\]

% ============================================================
\section{Recommended Final Course Outcomes}

Upon successful completion of the course, students will be able to:

\begin{enumerate}
    \item use algebraic, graphical, numerical, and symbolic representations
          of functions;
    \item construct and interpret elementary mathematical models;
    \item evaluate and interpret limits;
    \item determine continuity;
    \item calculate and interpret derivatives;
    \item apply differentiation rules accurately;
    \item analyze functions using first- and second-derivative information;
    \item solve optimization and related-rate problems;
    \item find antiderivatives;
    \item interpret and calculate definite integrals;
    \item use the Fundamental Theorem of Calculus;
    \item calculate cumulative change and average values;
    \item interpret calculus results in biological, health, behavioral, or
          social-science contexts;
    \item communicate mathematical reasoning clearly;
    \item recognize assumptions and limitations in mathematical models; and
    \item use calculus as a quantitative reasoning tool for scientific
          decision-making.
\end{enumerate}

% ============================================================
\section{Textbook-to-Course Crosswalk}

\begin{longtable}{p{2.3cm}p{4.0cm}p{6.3cm}p{2.2cm}}
\toprule
\textbf{Textbook} & \textbf{Course Module} & \textbf{Primary Role}
& \textbf{Status} \\
\midrule

Ch. 1
& Prerequisite Foundations
& Algebra, functions, trigonometry, graphing
& Review \\

Ch. 2
& Modeling Bridge
& Recurrence equations, sequences, exponential models
& Enrichment/Core bridge \\

Ch. 3
& Limits and Continuity
& Conceptual foundation for calculus
& Core \\

Ch. 4
& Differentiation
& Definition and mechanics of derivatives
& Core \\

5.1--5.6
& Applications of Differentiation
& Extrema, MVT, monotonicity, concavity, optimization
& Core \\

5.10
& Antiderivatives
& Preparation for integration
& Core \\

6.1
& Definite Integral
& Area, Riemann integral, accumulation
& Core \\

6.2
& Fundamental Theorem
& Differentiation/integration connection
& Core \\

6.3
& Applications of Integration
& Cumulative change and average values
& Core \\

7.1
& Substitution
& Integration technique
& Optional/Next course \\

7.2
& Integration by Parts
& Advanced integration
& Next course \\

8
& Differential Equations
& Continuous-time modeling
& Next course \\

9--12
& Advanced Quantitative Methods
& Linear algebra, multivariable calculus, systems, probability/statistics
& Subsequent courses \\

\bottomrule
\end{longtable}

% ============================================================
\section*{Final Design Principle}

The course should not be organized simply as

\[
\text{``limits} \rightarrow \text{derivatives} \rightarrow
\text{integrals.''}
\]

Instead, for Biology and Social Sciences students, the organizing principle
should be:

\[
\boxed{
\text{Phenomenon}
\rightarrow
\text{Model}
\rightarrow
\text{Mathematical Question}
\rightarrow
\text{Calculus Tool}
\rightarrow
\text{Interpretation}
}
\]

This preserves the central philosophy of the textbook: calculus is presented
as a rigorous mathematical framework for analyzing phenomena in the living
world, while students develop the ability to transfer mathematical concepts
to scientific situations.

\end{document}